% ID: SEL-D0010-BD-E08
% ESTADO: validada. S011; tres pasadas del mismo autor, sin segundo revisor.
\textbf{a) Desviación de las medias muestrales.}\par\nopagebreak[4]
Para una muestra aleatoria independiente de tamaño $n$ de una población normal con desviación conocida $\sigma$, la media sigue una normal con desviación $\sigma/\sqrt n$. Así,
\[
\boxed{\sigma_{\overline X}=\frac4{\sqrt{12}}=\frac2{\sqrt3}\approx1.154701.}
\]
La normalidad de la media es exacta bajo el modelo, aunque $n$ sea pequeño.

\textbf{b) Intervalo para la media desconocida.}\par\nopagebreak[4]
Los doce datos suman $127.2$, luego $\overline x=127.2/12=10.6$. Para confianza 97\,\%, $\alpha=0.03$ y $z_{0.985}\approx2.170090378$. Como la desviación poblacional es conocida, se usa el cuantil normal, no una desviación estimada a partir de la muestra:
\[
IC=\overline x\pm z_{0.985}\frac4{\sqrt{12}},\qquad
E\approx2.505804527.
\]
\[
\boxed{IC_{\text{97\%}}\approx(8.094195473,\ 13.105804527).}
\]
Con $z=2.17$ tabulado se obtiene prácticamente el mismo intervalo. La media poblacional es fija y desconocida; la confianza expresa la cobertura del método al repetir el muestreo.

\textbf{c) Error menor que 1.2.}\par\nopagebreak[4]
Se mantiene $\sigma=4$ y la confianza:
\[
z_{0.985}\frac4{\sqrt n}<1.2
\quad\Longleftrightarrow\quad
n>\left(\frac{4z_{0.985}}{1.2}\right)^2
\approx52.325469410.
\]
\boxed{n_{\min}=53}. El tamaño $52$ produce un error mayor que $1.2$, y $53$ ya da un error menor. La desigualdad es estricta; se toma el primer entero que la satisface. Con $2.17$ se mantiene el resultado $53$.
